
\section{Arquitectura del prototipo}
\label{sec:arquitectura}

Esta sección describe la organización del prototipo: primero la división en
paquetes y responsabilidades, después el flujo de datos del pipeline
determinista (RQ1 y RQ2) y el del subsistema de evaluación de los modelos de
lenguaje (RQ3), y por último las decisiones estructurales que lo sostienen.

\subsection{Visión general}
La organización del prototipo en paquetes responde a una idea sencilla:
cada paquete asume una responsabilidad y solo una. Se parte de las tres
responsabilidades clásicas de una herramienta de análisis y
refactorización, detección, transformación y medición; y se amplían
con dos piezas propias de este trabajo: un pipeline experimental que
orquesta los \emph{lotes} (cada lote es una ejecución del proceso sobre todo
el corpus) y un subsistema dedicado a la evaluación de los LLMs
(RQ3). Esta separación no es estética; es lo que permite cambiar una
parte sin arrastrar a las demás y, sobre todo, razonar sobre cada etapa de
forma aislada. Todas las operaciones sobre código se apoyan en el AST que
proporciona la biblioteca \textbf{JavaParser}, declarada como dependencia
en \texttt{pom.xml}. La \autoref{fig:paquetes} (diagrama de paquetes) y la
\autoref{fig:pipeline-det} (flujo de datos del pipeline determinista)
acompañan esta sección.

\subsection{Paquetes y responsabilidades}
El prototipo se organiza en paquetes con responsabilidades separadas. La
\autoref{tab:paquetes-arq} resume cada paquete y sus clases principales.

\begin{longtable}{@{}p{0.30\linewidth}p{0.62\linewidth}@{}}
\caption{Paquetes del prototipo y sus clases principales.}\label{tab:paquetes-arq}\\
\toprule
\textbf{Paquete} & \textbf{Responsabilidad y clases principales} \\
\midrule
\endhead
\texttt{es.tfm.refactoring} & Punto de entrada de línea de órdenes
(\texttt{Main}). \\
\texttt{\dots detection} & Análisis del AST y evaluación de P1--P4 (y del filtro opcional P5):
\texttt{NestedIfDetector}, \texttt{RefactoringOpportunity},
\texttt{DetectionResult}, \texttt{DiscardReason}, \texttt{DetectionMode},
\texttt{MethodCallAllowlist}. \\
\texttt{\dots transformation} & Aplicación de la transformación:
\texttt{NestedIfTransformer}. \\
\texttt{\dots analysis} & Cálculo proxy de complejidad cognitiva e impacto:
\texttt{CognitiveComplexityCalculator}, \texttt{RefactoringImpact},
\texttt{RefactoringImpactAnalyzer}, \texttt{CcAnalyzeCli}. \\
\texttt{\dots validation} & Comprobación de precondiciones y validez de la
transformación: \texttt{RefactoringValidator}, \texttt{ValidationResult}. \\
\texttt{\dots scanner} & Escaneo masivo de proyectos de código abierto para
construir el corpus real: \texttt{ProjectCorpusScanner},
\texttt{CorpusScanRunner}, \texttt{CorpusFileGenerator}, \texttt{ScanFinding}. \\
\texttt{\dots experiment} & Carga de corpus, ejecución por lotes,
exportación y contraste con SonarQube: \texttt{BatchRunner},
\texttt{Rq2BatchExecutor}, \texttt{ExperimentCase},
\texttt{ExperimentResult}, \texttt{ResultExporter},
\texttt{PilotCorpusLoader}, \texttt{RealDatasetLoader},
\texttt{SonarContrastEntry}, \texttt{SonarContrastReport}. \\
\texttt{\dots llm} & Protocolo, prompt, oráculo y campaña para RQ3:
\texttt{LlmExperimentProtocol}, \texttt{LlmPromptBuilder},
\texttt{LlmEvaluationSubset}, \texttt{LlmResponseValidator},
\texttt{LlmCampaignRunner}, \texttt{LlmCampaignReport},
\texttt{LlmCampaignExporter}, \texttt{LlmResponseProvider},
\texttt{LiveLlmResponseProvider}, \texttt{PrerecordedResponseProvider},
\texttt{LlmInvocationRecord}, \texttt{LlmEvaluationResult},
\texttt{LlmVerdict}, \texttt{LlmApiConfig}, \texttt{CampaignExecutor}. \\
\bottomrule
\end{longtable}

\noindent La tabla recoge las clases principales de cada paquete; se omiten
por brevedad algunos ejecutores y variantes auxiliares (p.\,ej.\
\texttt{Rq2ComparisonExecutor}, \texttt{Rq2ScannerBatchExecutor},
\texttt{TutorCorpusLoader}, \texttt{TrapCorpusLoader} en
\texttt{experiment}, o \texttt{Rq3PromptV2Executor} y
\texttt{Rq3TrapCampaignExecutor} en \texttt{llm}), que dan soporte a
experimentos secundarios y no alteran el flujo principal descrito.

\begin{figure}[H]
\centering
\begin{tikzpicture}[
  box/.style={draw, rounded corners, fill=blue!5, align=center, font=\small},
  orch/.style={box, fill=orange!12, text width=2.6cm, inner sep=4pt, font=\small\ttfamily},
  mainbox/.style={box, fill=orange!12, text width=5.8cm, inner sep=4pt, font=\small\ttfamily},
  core/.style={box, text width=9.4cm, inner sep=6pt},
  supp/.style={box, text width=2.6cm, inner sep=4pt, font=\small\ttfamily},
  arr/.style={-{Latex}, gray!70},
]
\node[mainbox] (main) at (4,0)      {Main / ExperimentOrchestrator};
\node[orch]    (exp)  at (0.5,-2.3) {experiment};
\node[orch]    (llm)  at (4,-2.3)   {llm};
\node[orch]    (scan) at (7.5,-2.3) {scanner};
\node[core]    (core) at (4,-4.7)
  {Núcleo de análisis y transformación\\[2pt]
   \ttfamily detection \quad transformation \quad analysis};
\node[supp]    (val)  at (0.5,-6.7) {validation};
% Entrada -> orquestación
\draw[arr] (main) -- (exp);
\draw[arr] (main) -- (llm);
\draw[arr] (main) -- (scan);
% Orquestación -> núcleo (verticales, sin cruces)
\draw[arr] (exp)  -- (exp  |- core.north);
\draw[arr] (llm)  -- (llm  |- core.north);
\draw[arr] (scan) -- (scan |- core.north);
% Núcleo -> validation
\draw[arr] (core.south -| val) -- (val);
\end{tikzpicture}
\caption{Diagrama de paquetes del prototipo. La capa de entrada
(\texttt{Main}) lanza los flujos de orquestación (\texttt{experiment} para
RQ1/RQ2, \texttt{llm} para RQ3, \texttt{scanner} para construir el corpus),
que se apoyan en el núcleo de análisis y transformación
(\texttt{detection}, \texttt{transformation}, \texttt{analysis}) y en
\texttt{validation}.}
\label{fig:paquetes}
\end{figure}

La \autoref{fig:paquetes} muestra la organización en paquetes; la
\autoref{fig:clases-nucleo} desciende al nivel de clases del núcleo
determinista, con sus métodos públicos y las enumeraciones de configuración
(\texttt{DetectionMode}) y de motivos de descarte (\texttt{DiscardReason}).

% TODO: FALTA la imagen figuras/img/fig-4.4.png (solo existe el fuente .mmd).
% Mientras no se genere, el PDF muestra la ruta como texto en lugar del diagrama.
% Generarla antes de compilar la versión final con:
%   mmdc -i figuras/src/fig-4.4-clases-nucleo.mmd -o figuras/img/fig-4.4.png -w 2600
\begin{figure}[H]
\centering
\includegraphics[width=\linewidth,height=0.85\textheight,keepaspectratio]{figuras/img/fig-4.4.png}
\caption{Diagrama de clases (UML) del núcleo determinista: detección
(\texttt{NestedIfDetector}), transformación (\texttt{NestedIfTransformer}),
medición (\texttt{CognitiveComplexityCalculator}), validación
(\texttt{RefactoringValidator}) y orquestación por lotes
(\texttt{BatchRunner}), con las clases de dominio \texttt{RefactoringOpportunity}
y \texttt{DetectionResult} y las enumeraciones \texttt{DetectionMode} (modos
\texttt{STRICT}/\texttt{RELAXED}) y \texttt{DiscardReason}. Las relaciones
siguen la notación UML: agregación (rombo) del orquestador sobre los servicios
del núcleo, asociación con el modo de detección y dependencias de creación y
uso.}
\label{fig:clases-nucleo}
\end{figure}

\subsection{Flujo de datos del pipeline determinista (RQ1, RQ2)}

El flujo extremo a extremo, implementado por \texttt{BatchRunner.processCase}, se organiza en las siguientes etapas:

\begin{enumerate}
    \item \textbf{Carga del corpus.}  
    Los casos se cargan desde el \emph{classpath} y se construyen como instancias de \texttt{ExperimentCase}; cada caso es un método Java acompañado de sus metadatos (procedencia, complejidad esperada, etc.).

    \item \textbf{Parseo a AST.}  
    El código fuente se parsea mediante \lstinline|StaticJavaParser.parse(...)| y se selecciona el primer \lstinline|MethodDeclaration|. Si el parseo falla, el caso se marca como inelegible con un motivo explícito.

    \item \textbf{Detección con motivos.}  
    Se ejecuta \lstinline|detectWithReasons| para identificar oportunidades y registrar los \texttt{DiscardReason} asociados en \lstinline|discardCategories|.

    \item \textbf{Medición \emph{before}.}  
    Se calcula la complejidad cognitiva inicial del método antes de aplicar ninguna transformación.

    \item \textbf{Aplicación iterativa de la transformación.}  
    Sobre un clon del AST se aplica la política \emph{one-at-a-time} (véase \autoref{sec:one-at-a-time}), re-detectando tras cada cambio hasta alcanzar un punto fijo o el límite de seguridad.

    \item \textbf{Medición \emph{after}.}  
    Una vez finalizado el proceso iterativo, se vuelve a calcular la complejidad cognitiva y se obtiene $\Delta = \text{after} - \text{before}$.

    \item \textbf{Construcción del resultado experimental.}  
    Con la información anterior se construye una instancia inmutable de \texttt{ExperimentResult}.

    \item \textbf{Exportación de resultados.}  
    Finalmente, los resultados se exportan a JSON (incluyendo código fuente) y a CSV (formato tabular) (\autoref{fig:pipeline-det}).
\end{enumerate}

\begin{figure}[H]
\centering
\begin{tikzpicture}[
  term/.style={draw, rounded rectangle, fill=green!12, align=center,
               text width=3cm, inner sep=4pt, font=\small},
  endbad/.style={draw, rounded rectangle, fill=red!10, align=center,
                 text width=3cm, inner sep=4pt, font=\small},
  proc/.style={draw, fill=blue!5, align=center,
               text width=3.8cm, inner sep=5pt, font=\small},
  io/.style={draw, trapezium, trapezium left angle=70,
             trapezium right angle=110, trapezium stretches=true,
             fill=gray!15, align=center, text width=3.6cm, inner sep=4pt,
             font=\small},
  dec/.style={draw, diamond, aspect=2.2, fill=orange!12, align=center,
              text width=2.2cm, inner sep=1pt, font=\scriptsize},
  arr/.style={-{Latex}},
  lbl/.style={font=\scriptsize, inner sep=1.5pt},
]
% Eje principal
\node[term] (start) at (0,0)      {Inicio: \texttt{ExperimentCase}};
\node[dec]  (dpar)  at (0,-2.3)   {\textquestiondown{}AST válido y con método?};
\node[proc] (db)    at (0,-4.6)   {Detección con motivos\\+ CC \emph{before} (original)};
\node[proc] (loop)  at (0,-6.2)   {Detectar oportunidades\\(sobre el clon)};
\node[dec]  (dloop) at (0,-8.5)   {\textquestiondown{}oportunidad y pasadas $<$ \texttt{MAX}?};
\node[proc] (ar)    at (0,-10.8)  {CC \emph{after}, $\Delta$\\y \texttt{ExperimentResult}};
\node[io]   (exp)   at (0,-12.6)  {Exportar JSON + CSV};
\node[term] (fin)   at (0,-14.2)  {Fin};
% Ramas laterales
\node[endbad] (inel) at (4.8,-2.3) {Inelegible\\(\texttt{CC before = -1})};
\node[proc] (apply)  at (4.8,-8.5) {Aplicar la primera oportunidad};
% Flechas del eje
\draw[arr] (start) -- (dpar);
\draw[arr] (dpar)  -- node[lbl,left]{sí} (db);
\draw[arr] (db)    -- (loop);
\draw[arr] (loop)  -- (dloop);
\draw[arr] (dloop) -- node[lbl,left]{no} (ar);
\draw[arr] (ar)    -- (exp);
\draw[arr] (exp)   -- (fin);
% Ramas
\draw[arr] (dpar)  -- node[lbl,above]{no} (inel);
\draw[arr] (dloop) -- node[lbl,above]{sí} (apply);
\draw[arr] (apply) |- node[lbl,above]{re-detectar} (loop.east);
\end{tikzpicture}
\caption{Flujo de datos del pipeline determinista
(\texttt{BatchRunner.processCase}). Se hacen explícitos la rama de
inelegibilidad (fallo de parseo o ausencia de método $\rightarrow$
\texttt{CC before = -1}; el resultado inelegible también se registra y
exporta) y el bucle de la política \emph{one-at-a-time} (detectar
$\rightarrow$ aplicar la primera oportunidad $\rightarrow$ re-detectar,
hasta punto fijo o \texttt{MAX\_PASSES}). La CC \emph{before} se mide sobre
el método original y la \emph{after} sobre el clon transformado.}
\label{fig:pipeline-det}
\end{figure}

\subsection{Flujo del subsistema LLM (RQ3)}

Para RQ3 se añade un flujo paralelo coordinado por \texttt{CampaignExecutor} y \texttt{LlmCampaignRunner}, que se organiza en las siguientes etapas:

\begin{enumerate}
    \item \textbf{Carga del subset de evaluación.}  
    Se carga el subconjunto RQ3 definido en \texttt{LlmEvaluationSubset}, junto con sus correspondientes baselines deterministas.

    \item \textbf{Construcción del prompt.}  
    El mensaje enviado al modelo se genera mediante \texttt{LlmPromptBuilder}, utilizando la versión fijada del protocolo (\lstinline|PROMPT_VERSION = "v1.0"|).

    \item \textbf{Invocación al modelo.}  
    La ejecución del modelo se realiza a través de la interfaz \texttt{LlmResponseProvider}, con implementaciones específicas como \texttt{LiveLlmResponseProvider} y \texttt{PrerecordedResponseProvider}.

    \item \textbf{Validación mediante el oráculo.}  
    Cada respuesta se procesa con \texttt{LlmResponseValidator}, que aplica el oráculo experimental descrito en la \autoref{fig:oraculo}.

    \item \textbf{Agregación y exportación de resultados.}  
    Finalmente, los resultados se consolidan y exportan en artefactos.
\end{enumerate}

\subsection{Decisiones estructurales}
\begin{itemize}
  \item \textbf{Inmutabilidad de los resultados.} \texttt{ExperimentResult}
  y \texttt{LlmEvaluationResult} se construyen una sola vez y se exportan
  sin mutación posterior, lo que facilita la trazabilidad.
  \item \textbf{Provider intercambiable para LLM.}
  \texttt{LlmResponseProvider} permite ejecutar la misma campaña con
  respuestas reales o grabadas, desacoplando la lógica experimental del
  coste y la variabilidad de las llamadas remotas.
  \item \textbf{Exportadores múltiples.} CSV, JSON y Markdown cubren usos
  distintos (tabla agregada, evidencia completa, lectura humana).
  \item \textbf{Sin acoplamiento a SonarQube en el pipeline determinista.}
  El prototipo \emph{no invoca} SonarQube en el flujo principal; la
  validación cruzada se realiza de forma independiente del pipeline, como un
  contraste externo (\autoref{sec:sonar-val}).
\end{itemize}
